\begin{abcliste}
\abc Die Tabelle \emph{Dichte von Wasser und von Quecksilber} (FoTaBe S.~192) zeigt die grösste Dichte, \rMxO, bei $\vartheta = \result{\tMxrP{0}{1}}$.

\abc Die Masse des Wassers bleibt beim Erhitzen gleich: $m = \rho_{20}\,V_0 = \rho_{90}\,V_{90}$, mit $\rho_{20} = \rKaO$ und $\rho_{90} = \rKeO$. Das Wasser dehnt sich auf $V_{90} = V_0\,\rho_{20}/\rho_{90}$ aus. Wenn wir die Ausdehnung des Kanisters vernachlässigen, läuft die Differenz über:
\begin{align}
 \Delta V &= V_{90}-V_0 = \dVKF \\
   &= \VKaO\cdot\left(\frac{\rKaO}{\rKeO}-1\right) \\
   &= \dVKl \approx \result{\dVKlP{0}{2}}
\end{align}
Der Kanister dehnt sich beim Erhitzen ebenfalls etwas aus; tatsächlich läuft deshalb etwas weniger Wasser über.
\end{abcliste}
